\section*{Introduction}

Modern volume electron microscopy (vEM) routinely produces datasets in the 1--100~TB range, with focused-ion-beam scanning electron microscopy (FIB-SEM) volumes from the OpenOrganelle collection \cite{Xu2021} as a representative example. Deep-learning segmentation models are now the standard tool for extracting biological structure from these volumes \cite{Heinrich2021,Funke2018}, and trained models exist for a growing list of organelles. The bottleneck has shifted from \emph{producing} predictions to \emph{deploying} them: a researcher who wants to look at predicted mitochondria across a 10~TB dataset, compare two models on the same region, or iterate on a training run without re-exporting the entire output volume, runs into the same recurring pain points.

Three problems dominate. First, full volumes do not fit on a single GPU, so inference must be chunked, but the chunking, padding, and re-assembly logic is repeatedly reimplemented in ad-hoc scripts. Second, results need to be explored interactively at multiple scales --- the natural format is something like Neuroglancer's precomputed multiscale Zarr \cite{Neuroglancer}, but populating that format up front means hours of compute before the first qualitative look. Third, scaling out to a cluster, retrying failed blocks, recovering from preemption, and combining predictions from several models all demand a control plane that the model itself does not provide.

Existing tools address pieces of this. The BioImage.io zoo \cite{BioImageIO} standardises model packaging and provides ONNX-based serving, but does not integrate with chunked viewers. DaCapo \cite{Patton2024} and Funkelab pipelines support training and offline batch inference but not interactive exploration. Neuroglancer's precomputed format requires converting predictions to disk before viewing. None of these tools, individually, lets a user point a model at a multi-terabyte dataset and immediately scroll through predictions in a browser.

We present \cellmapflow{}, an open-source framework that closes this gap by treating inference as a service. Predictions are exposed as a virtual Zarr volume over HTTP, computed on demand by GPU workers as Neuroglancer requests chunks. The same execution engine, behind the same configuration, also drives whole-volume batch inference on a cluster using a Daisy-based block scheduler \cite{DaisyFunke}. The system accepts models from multiple ecosystems (PyTorch TorchScript, DaCapo, BioImage.io, HuggingFace Hub, lab-internal formats), reads volumetric data from Zarr, N5, OME-Zarr, S3, and Google Cloud precomputed via TensorStore, composes user-supplied normalizers and post-processors into YAML-defined pipelines, and supports multi-model fusion at inference time. \cellmapflow{} is in active production use within the CellMap project at Janelia for segmentation across the OpenOrganelle dataset collection.

In the remainder of this paper we describe the system architecture (Section~\ref{sec:architecture}), the post-processing and multi-model fusion layer (Section~\ref{sec:postprocessing}), demonstrate representative use cases (Section~\ref{sec:usecases}), report a benchmark suite covering interactive latency, blockwise throughput, and cluster scaling (Section~\ref{sec:benchmarks}), and discuss limitations and future directions (Section~\ref{sec:discussion}).
